\input{../header}
\usepackage{pgfplots}
\rhead{Your name: \rule{8cm}{0.15mm}}

\begin{document}

\section{Learning target IN3 and IN5, version 1} % checkit 15

Evaluate a definite integral using the Fundamental Theorem of Calculus.

\hrulefill


Explain how to compute the exact value of each of the following definite integrals using the Fundamental Theorem of Calculus. 

\textbf{Leave all answers in exact form, with no decimal approximations. }

\begin{enumerate}[(a)]
\item
 \(\displaystyle \int_{ x=4 }^{ 6 } \left( \frac{6}{x} \right) dx\) 

 \vfill

\item
 \(\int_{x= \frac{5}{4} \, \pi }^{ \frac{4}{3} \, \pi } \left( -2 \, \sec^2\left(x\right) \right) dx\)

\vfill

\item
 \(\int_{x= -2 }^{ -1 } \left( 3 \, x^{3} + 10 \, x - 1 \right) dx\)
\vfill
\end{enumerate}
%%%%%%%%%%%%%%%%%%%%%%%%%%%%%%%%%%%%%%%%%%%%%%%%%%%%%%%%%
\pagebreak

\section{Learning target IN3 and IN5, version 2} 

Evaluate a definite integral using the Fundamental Theorem of Calculus.

\hrulefill


Explain how to compute the exact value of each of the following definite integrals using the Fundamental Theorem of Calculus. 

\textbf{Leave all answers in exact form, with no decimal approximations. }

\begin{enumerate}[(a)]
\item
\(\displaystyle \int_{ x=1 }^{ 5 } \left( \frac{2}{x} + 2^x\right) dx\) 
\vfill
\item 
\(\displaystyle \int_{x=\pi/3}^{3\pi/4} \left( 7\cos(x) \right) dx\)
\vfill
\item
\(\displaystyle \int_{x= -3 }^{ 1 } \left( 9x^3 -9x^2 + 2\right) dx\)
\vfill

\end{enumerate}
%%%%%%%%%%%%%%%%%%%%%%%%%%%%%%%%%%%%%%%%%%%%%%%%%%%%%%%%%
\pagebreak

\section{Learning target IN3 and IN5, version 3} 

Evaluate a definite integral using the Fundamental Theorem of Calculus.

\hrulefill


Explain how to compute the exact value of each of the following definite integrals using the Fundamental Theorem of Calculus. 

\textbf{Leave all answers in exact form, with no decimal approximations. }

\begin{enumerate}[(a)]
\item
\(\displaystyle \int_{x= -2 }^{ 4 } \left( -3x^2 + 4x - 12\right) dx\)
\vfill
\item 
\(\displaystyle \int_{ x=2 }^{ 3 } \left( e^x - \dfrac{3}{x}\right) dx\) 
\vfill
\item
\(\displaystyle \int_{x=\pi/6}^{2\pi/3} \left( 4\sin(x) \right) dx\)
\vfill

\end{enumerate}
%%%%%%%%%%%%%%%%%%%%%%%%%%%%%%%%%%%%%%%%%%%%%%%%%%%%%%%%%
\pagebreak
\section{Learning target IN3 and IN5, version 4} 

Evaluate a definite integral using the Fundamental Theorem of Calculus.

\hrulefill


Explain how to compute the exact value of each of the following definite integrals using the Fundamental Theorem of Calculus. 

\textbf{Leave all answers in exact form, with no decimal approximations. }

\begin{enumerate}[(a)]
\item
\(\displaystyle \int_{x= 0 }^{ \pi } \left( 7\sin(x) - 2x^2\right) dx\)
\vfill
\item 
\(\displaystyle \int_{ x=2 }^{ 10 } \left( \dfrac{3}{x} + 7^x\right) dx\) 
\vfill
\item
\(\displaystyle \int_{x=-1}^{3} \left(9x^4-12x^2+7x-3 \right) dx\)
\vfill

\end{enumerate}
%%%%%%%%%%%%%%%%%%%%%%%%%%%%%%%%%%%%%%%%%%%%%%%%%%%%%%%%%
\pagebreak
%%%%%%%%%%%%%%%%%%%%%%%%%%%%%%%%%%%%%%%%%%%%%%%%%%%%%%%%%
\end{document}